\documentclass[11pt,a4paper]{article}

\usepackage[utf8]{inputenc}
\usepackage[T1]{fontenc}
\usepackage[brazil]{babel}
\usepackage[margin=2.5cm]{geometry}
\usepackage{amsmath}
\usepackage{graphicx}
\usepackage{booktabs}
\usepackage{siunitx}
\usepackage{float}
\usepackage{listings}
\usepackage{xcolor}
\usepackage[hidelinks]{hyperref}

\sisetup{output-decimal-marker={,}}
\graphicspath{{plots/}}

\lstset{
  basicstyle=\ttfamily\scriptsize,
  keywordstyle=\color{blue!60!black},
  commentstyle=\color{green!40!black},
  stringstyle=\color{red!60!black},
  numbers=left,
  numberstyle=\tiny\color{gray},
  breaklines=true,
  frame=single,
  language=Matlab,
  showstringspaces=false,
  captionpos=b
}

\title{Aplicações de Processamento Digital de Sinais -- 4456S-04\\[2mm]
\large Prática 7: Compressão de Imagem}
\author{Antônio Vinicius Corrêa dos Santos}
\date{13/10/2026}

\begin{document}

\maketitle

Imagem: \texttt{lena\_c.jpg}, $256\times256$, convertida para escala de cinza com
\texttt{rgb2gray} e mantida em $[0,255]$. Todas as figuras são geradas por
\texttt{main.m} e gravadas em \texttt{plots/}.

\section{Espectro da DCT 2D}

\begin{figure}[H]
\centering
\includegraphics[width=0.49\textwidth]{dct_spectrum.jpg}
\hfill
\includegraphics[width=0.49\textwidth]{dct_spectrum_zoom.jpg}
\caption{Espectro de magnitude da DCT, $\log(1+|X(u,v)|)$: matriz completa e canto
$32\times32$ de baixa frequência.}
\end{figure}

Os pontos mais brilhantes estão no \textbf{canto superior esquerdo} da matriz, em torno
de $(u,v)=(1,1)$, e a amplitude cai ao longo das diagonais em direção ao canto inferior
direito. O coeficiente $(1,1)$ é a componente DC, proporcional ao brilho médio da
imagem; os vizinhos são as baixas frequências espaciais, que descrevem as variações
lentas de intensidade --- regiões homogêneas, sombreamento e a forma geral dos objetos.
Os coeficientes distantes do canto representam as variações rápidas (bordas finas,
textura do cabelo e do chapéu, ruído) e têm amplitude muito menor.

Diferente da DFT, a DCT não precisa de \texttt{fftshift}: ela não tem frequências
negativas, então a origem já é o canto, e como a extensão implícita da imagem é par,
não aparecem as faixas causadas pela descontinuidade entre bordas opostas vistas na
Prática 6.

\section{Máscara zonal triangular}

A máscara $A$ vale $1$ no triângulo superior esquerdo de lado $j$, ou seja, nas posições
com $u+v \le j+1$, e $0$ nas demais:
\par\noindent\begin{minipage}{\linewidth}
\begin{lstlisting}[numbers=none,frame=none]
function A = zonal_mask(j, N)
    A = zeros(N);
    A(1:j,1:j) = fliplr(triu(ones(j)));
end
\end{lstlisting}
\end{minipage}
\texttt{triu} gera o triângulo acima da diagonal principal e \texttt{fliplr} o espelha
para que a hipotenusa fique na antidiagonal. O número de coeficientes mantidos é
$j(j+1)/2$.

\begin{figure}[H]
\centering
\includegraphics[width=0.9\textwidth]{masks_and_error.jpg}
\caption{Máscara para $j=8$ (canto $16\times16$), máscara para $j=128$ e erro absoluto
da reconstrução com $j=128$.}
\end{figure}

\section{Reconstrução com coeficientes truncados}

Para cada $j\in\{1,2,4,8,16,32,64,128\}$: $Y = \text{idct2}(X\odot A_j)$. As
reconstruções em tamanho original estão em \texttt{plots/recon\_j-*.jpg}.

\begin{figure}[H]
\centering
\includegraphics[width=0.55\textwidth]{recon_grid.jpg}
\caption{Original e reconstruções para os oito valores de $j$.}
\end{figure}

\begin{figure}[H]
\centering
\includegraphics[width=0.3\textwidth]{recon_j-016.jpg}
\hfill
\includegraphics[width=0.3\textwidth]{recon_j-032.jpg}
\hfill
\includegraphics[width=0.3\textwidth]{recon_j-128.jpg}
\caption{Reconstruções para $j=16$, $32$ e $128$, a partir dos arquivos em resolução original ($256\times256$).}
\end{figure}

\begin{table}[H]
\centering
\caption{Coeficientes mantidos, fração da energia retida
$\sum|X\odot A|^2/\sum|X|^2$ e erro em relação à original.}
\begin{tabular}{rrrccc}
\toprule
$j$ & \textbf{Coef.} & \textbf{Mantidos (\%)} & \textbf{Energia} & \textbf{RMSE} &
\textbf{PSNR (dB)} \\
\midrule
$1$   & $1$    & $0{,}0015$  & $0{,}8708$ & $47{,}77$ & $14{,}55$ \\
$2$   & $3$    & $0{,}0046$  & $0{,}8794$ & $46{,}15$ & $14{,}85$ \\
$4$   & $10$   & $0{,}0153$  & $0{,}8913$ & $43{,}83$ & $15{,}30$ \\
$8$   & $36$   & $0{,}0549$  & $0{,}9202$ & $37{,}55$ & $16{,}64$ \\
$16$  & $136$  & $0{,}2075$  & $0{,}9527$ & $28{,}90$ & $18{,}91$ \\
$32$  & $528$  & $0{,}8057$  & $0{,}9740$ & $21{,}44$ & $21{,}51$ \\
$64$  & $2080$ & $3{,}1738$  & $0{,}9871$ & $15{,}07$ & $24{,}57$ \\
$128$ & $8256$ & $12{,}5977$ & $0{,}9946$ & $9{,}74$  & $28{,}36$ \\
\bottomrule
\end{tabular}
\end{table}

Observações:

\begin{itemize}\itemsep0pt
  \item \textbf{$j=1$ a $4$.} Com só o termo DC a imagem é um cinza uniforme no nível
        médio; até $j=4$ aparecem apenas gradientes suaves. Ainda assim o termo DC
        sozinho carrega \SI{87}{\percent} da energia --- energia não é o mesmo que
        informação visual.
  \item \textbf{$j=8$ e $16$.} Surgem manchas claras e escuras na posição do rosto e do
        chapéu, mas sem forma definida.
  \item \textbf{A partir de $j=32$} as características principais tornam-se
        discerníveis: o contorno do chapéu, o rosto, o ombro e a moldura ao fundo são
        reconhecíveis, embora borrados. Nesse ponto são mantidos
        $32\cdot33/2 = 528$ coeficientes, isto é,
        $528/65536 \approx \SI{0,81}{\percent}$ do total.
  \item \textbf{$j=64$} já mostra os olhos, a boca e a pena do chapéu com
        \SI{3,2}{\percent} dos coeficientes.
  \item O borrado das reconstruções tem um leve \emph{ringing} junto às bordas fortes,
        efeito do corte abrupto da máscara (o mesmo fenômeno do passa-baixas ideal da
        Prática 6).
\end{itemize}

\section{Comparação entre $j=128$ e a original}

Com $j=128$ são descartados \SI{87,4}{\percent} dos coeficientes e, à distância normal
de visualização, a imagem é praticamente indistinguível da original: PSNR de
\SI{28,4}{dB} e \SI{99,46}{\percent} da energia retida. A perda existe, mas é sutil:
uma leve suavização da textura do cabelo e da pena, e o mapa de erro mostra que ela
está concentrada exatamente nessas bordas e texturas finas, enquanto as regiões
homogêneas (pele, fundo) têm erro quase nulo.

A justificativa é a \textbf{compactação de energia} da DCT. Pixels vizinhos de uma
imagem natural são fortemente correlacionados, e a DCT se aproxima da transformada
ótima (Karhunen--Loève) para esse tipo de sinal: ela concentra quase toda a energia em
poucos coeficientes de baixa frequência. Os \SI{87}{\percent} descartados respondem por
apenas \SI{0,54}{\percent} da energia, e correspondem a detalhes de alta frequência aos
quais o olho humano é pouco sensível. É esse o princípio explorado pelo JPEG, que aplica
a DCT em blocos $8\times8$ e quantiza com mais força os coeficientes de alta frequência.

\section{DCT em blocos com \texttt{blockproc}}

A imagem foi dividida em blocos $B\times B$ e a DCT aplicada em cada bloco
independentemente, com \texttt{funcao\_dct = @(b) dct2(b.data)} e, na volta,
\texttt{funcao\_idct = @(b) idct2(b.data)}. No Octave o \texttt{blockproc} entrega o
bloco como matriz, e não como \emph{struct} com campo \texttt{data}; a função auxiliar
\texttt{bdata} aceita os dois formatos, de modo que o script roda igual no MATLAB.

\subsection*{Pixels da imagem original e transformada ($B=8$)}

O primeiro bloco $8\times8$ (canto superior esquerdo, região do fundo) e sua DCT:
\[
x = \begin{bmatrix}
161 & 161 & 162 & 163 & 164 & 163 & 160 & 158\\
162 & 161 & 162 & 162 & 162 & 161 & 158 & 156\\
161 & 161 & 161 & 161 & 161 & 159 & 157 & 154\\
159 & 159 & 159 & 159 & 159 & 158 & 156 & 154\\
157 & 156 & 156 & 157 & 158 & 158 & 156 & 155\\
156 & 155 & 155 & 156 & 157 & 158 & 157 & 156\\
156 & 156 & 155 & 156 & 157 & 157 & 156 & 155\\
158 & 157 & 156 & 156 & 156 & 156 & 155 & 154
\end{bmatrix}
\qquad
X = \begin{bmatrix}
1264 & 6 & -6 & 5 & 0 & 0 & 0 & 0\\
16 & 4 & -4 & 0 & 0 & 0 & 0 & 0\\
4 & 0 & 0 & 0 & 0 & 0 & 0 & 0\\
0 & -6 & 0 & 1 & 0 & 0 & 0 & 0\\
0 & 0 & 0 & 0 & 0 & 0 & 0 & 0\\
0 & 0 & -1 & 0 & 0 & 0 & -1 & 0\\
0 & 0 & 0 & 0 & 0 & 0 & 0 & 0\\
0 & 0 & 0 & 0 & 0 & 0 & 0 & 0
\end{bmatrix}
\]
(valores arredondados). Os pixels variam pouco em torno de $158$, e a DCT reflete isso:
o termo DC vale $1264 = 8\cdot\bar{x}$ ($\bar{x}\approx158$) e quase todos os demais
coeficientes são próximos de zero. A compactação de energia observada na imagem inteira
se repete dentro de cada bloco. Na imagem transformada completa
(Figura~\ref{fig:block}, à esquerda) cada bloco tem seu ponto brilhante no canto
superior esquerdo, formando uma grade com a mesma periodicidade dos blocos.

\subsection*{Zerando coeficientes}

Em cada bloco foram mantidos apenas os coeficientes do triângulo $u+v\le B/2+1$ (a
mesma máscara zonal da seção anterior, com $j=B/2$) e o resto foi zerado. No primeiro
bloco isso remove só o $-6$ e o $1$ da quarta linha e os $-1$ da sexta; a reconstrução
fica
\[
\hat{x} = \begin{bmatrix}
162 & 162 & 163 & 163 & 164 & 162 & 159 & 157\\
161 & 161 & 162 & 162 & 163 & 161 & 159 & 156\\
160 & 160 & 160 & 161 & 161 & 160 & 158 & 155\\
159 & 158 & 158 & 159 & 159 & 159 & 157 & 155\\
157 & 157 & 157 & 157 & 158 & 157 & 156 & 154\\
157 & 156 & 156 & 156 & 157 & 157 & 156 & 154\\
157 & 156 & 155 & 156 & 157 & 157 & 156 & 155\\
157 & 156 & 155 & 156 & 157 & 157 & 156 & 155
\end{bmatrix},
\]
com erro de no máximo $\pm2$ níveis de cinza por pixel. Zerar um coeficiente da DCT não
apaga um pixel: cada coeficiente é uma função-base cosseno que cobre o bloco inteiro,
então o erro se espalha de forma suave por todos os $64$ pixels do bloco, e nenhum
pixel fora dele é afetado.

Sem zerar nada, a IDCT recupera a imagem original até o erro numérico
($\max|\hat{x}-x| < 10^{-12}$), confirmando que a DCT em blocos é inversível.

\subsection*{Variação do tamanho do bloco}

Os pixels foram verificados novamente na mesma janela $8\times8$ do canto superior
esquerdo, agora para $B=4$, $16$ e $32$. Na imagem transformada:
\[
X_{B=4} = \begin{bmatrix}
644 & -1 & 1 & 0 & 635 & 9 & -1 & 0\\
4 & -1 & 0 & 0 & 7 & 1 & 0 & 0\\
-2 & -1 & 0 & 0 & 1 & -1 & 0 & 0\\
1 & 0 & 0 & 0 & 0 & 1 & 0 & 1\\
625 & 1 & 2 & 0 & 625 & 3 & -1 & -1\\
0 & -1 & 0 & 0 & 2 & 0 & 0 & 0\\
2 & 1 & 0 & 0 & -1 & 1 & 0 & 0\\
0 & 0 & 0 & 0 & 0 & 1 & 0 & 0
\end{bmatrix}
\]
\[
X_{B=32} = \begin{bmatrix}
4971 & 138 & -144 & 239 & -178 & 62 & -105 & 89\\
40 & -30 & -60 & -31 & 15 & 53 & 27 & -29\\
-36 & -40 & 93 & 90 & -43 & -45 & -7 & 2\\
40 & 32 & -17 & -55 & -18 & 48 & 25 & -4\\
0 & -10 & 12 & 19 & 25 & 1 & -33 & -7\\
10 & 2 & 5 & -20 & -8 & -4 & 6 & 6\\
-1 & -1 & 7 & 8 & 12 & -6 & -4 & 0\\
14 & 8 & -5 & -13 & -2 & -3 & 6 & 0
\end{bmatrix}
\]

\begin{itemize}\itemsep0pt
  \item Com $B=4$ a janela contém quatro blocos inteiros, e aparecem quatro termos DC
        ($644$, $635$, $625$, $625$), cada um igual a $4$ vezes a média do seu bloco
        ($\approx161$, $159$, $156$, $156$). Os outros coeficientes são quase nulos:
        num bloco tão pequeno de fundo quase não há variação.
  \item Com $B=16$ e $32$ a janela mostra só o canto de baixa frequência de um único
        bloco. O DC cresce com o bloco ($2516 = 16\cdot157{,}3$ e
        $4971 = 32\cdot155{,}3$), e os coeficientes AC deixam de ser desprezíveis
        (até $\approx240$ para $B=32$), porque o bloco já abrange o fundo, a moldura e
        parte do chapéu.
  \item Com $B=32$ a janela inteira está dentro do triângulo mantido ($u+v\le17$), então
        nenhum dos coeficientes mostrados é zerado; o que se perde fica fora dela.
\end{itemize}

Na imagem reconstruída após zerar os coeficientes, os pixels da janela continuam
próximos da original para todos os tamanhos. O erro máximo nessa janela foi
$1{,}42$, $1{,}55$, $1{,}88$ e $2{,}34$ níveis de cinza para $B=4$, $8$, $16$ e $32$.
Ele cresce um pouco com $B$ porque o erro de cada bloco se espalha por mais pixels,
mas continua imperceptível nessa região de fundo.

\begin{figure}[H]
\centering
\includegraphics[width=0.8\textwidth]{block_04.jpg}\\
\includegraphics[width=0.8\textwidth]{block_08.jpg}\\
\includegraphics[width=0.8\textwidth]{block_16.jpg}\\
\includegraphics[width=0.8\textwidth]{block_32.jpg}
\caption{Para $B=4,8,16,32$: DCT em blocos ($\log(1+|X|)$), reconstrução após zerar
os coeficientes fora do triângulo $j=B/2$ e \emph{zoom} na região do olho.}
\label{fig:block}
\end{figure}

\begin{table}[H]
\centering
\caption{Reconstrução com coeficientes zerados, por tamanho de bloco.}
\begin{tabular}{rrcc}
\toprule
$B$ & \textbf{Mantidos (\%)} & \textbf{RMSE} & \textbf{PSNR (dB)} \\
\midrule
$4$  & $18{,}75$ & $9{,}94$ & $28{,}18$ \\
$8$  & $15{,}62$ & $9{,}89$ & $28{,}23$ \\
$16$ & $14{,}06$ & $9{,}84$ & $28{,}27$ \\
$32$ & $13{,}28$ & $9{,}81$ & $28{,}29$ \\
\bottomrule
\end{tabular}
\end{table}

\begin{itemize}\itemsep0pt
  \item Na imagem transformada a grade acompanha $B$: com blocos pequenos ela parece
        uma versão em mosaico da própria imagem (cada termo DC é a média do bloco);
        com blocos grandes cada bloco mostra um espectro de verdade, com a energia
        caindo do canto para fora.
  \item Blocos maiores mantêm uma fração menor dos coeficientes e ainda assim dão erro
        um pouco menor: há mais pixels correlacionados por bloco, então a compactação
        de energia é melhor.
  \item O tipo de artefato muda com $B$. Com blocos pequenos ($4$, $8$) aparece o
        \emph{efeito de bloco}: cada bloco é aproximado sozinho, e as bordas entre
        blocos vizinhos ficam visíveis como degraus no \emph{zoom}. Com blocos grandes
        ($16$, $32$) as fronteiras somem, mas volta o \emph{ringing} junto às bordas
        fortes, como na DCT da imagem inteira.
  \item O custo computacional também cresce com $B$, e blocos grandes misturam
        regiões de conteúdo diferente. O JPEG adota $8\times8$ como compromisso entre
        compactação, custo e localização dos artefatos.
\end{itemize}

\section*{Script (Octave)}

\lstinputlisting[caption={\texttt{main.m}}]{main.m}

\end{document}
